\documentclass[10pt,a4paper,reqno]{amsart}
\usepackage{amsthm,amssymb,amstext,graphicx,comment}
\usepackage[foot]{amsaddr}
\usepackage[a4paper, margin=2.5cm]{geometry}
\usepackage{enumerate}
\usepackage{amscd}
\usepackage{mathtools}
\usepackage[bibliography= common]{apxproof}
\usepackage[linesnumbered,ruled,vlined]{algorithm2e}
\SetKwInput{KwInput}{Input}                % Set the Input
\SetKwInput{KwOutput}{Output}              % set the Output
\newcommand{\KwContinue}{\bf continue}
\usepackage{hyperref}
\usepackage{booktabs}
\usepackage{float}
\usepackage[table]{xcolor}% http://ctan.org/pkg/xcolor
\usepackage{mathabx}
\usepackage{bm}
\newcommand{\XiI}{\Xi_1^{11}}
\newcommand{\F}{\mathcal{F}}
\newcommand{\XiII}{\Xi_1^{12}}
\newcommand{\XiIII}{\Xi_1^{21}}
\newcommand{\XiIV}{\Xi_1^{22}}
\usepackage{bbm}
\usepackage{subcaption}
\usepackage{array}
\usepackage{empheq}
% \newcommand{\Hrow}{1.2}        % vertical spacing between rows
% \newcommand{\W}{12}            % horizontal axis length
% \newcommand{\h}{0.45}          % rectangle height

% \tikzset{
%   monthline/.style={thick},
%   tick/.style      ={thick},
%   tradable/.style  ={draw=black, fill=gray!30},
%   held/.style      ={draw=black, fill=blue!30},
%   lbl/.style       ={font=\small}
% }

% \usetikzlibrary{patterns.meta}

% \tikzset{
%     hash/.style={
%         pattern={Lines[angle=45,distance=4pt,line width=.4pt]},
%         pattern color=black!50
%     },
%     bluebox/.style={fill=blue!40},
%     monthline/.style={thick,->},
%     tick/.style={line width=.4pt},
%     lbl/.style={font=\small}
% }

\usepackage{textcomp}
\usepackage{multirow}
\hypersetup{
  colorlinks=true,
  pdfauthor={© Sven Karbach},
  pdfsubject={},
  pdfkeywords={}
}
\makeatletter
\newsavebox{\@brx}
\newcommand{\llangle}[1][]{\savebox{\@brx}{\(\m@th{#1\langle}\)}%
  \mathopen{\copy\@brx\kern-0.5\wd\@brx\usebox{\@brx}}}
\newcommand{\rrangle}[1][]{\savebox{\@brx}{\(\m@th{#1\rangle}\)}%
  \mathclose{\copy\@brx\kern-0.5\wd\@brx\usebox{\@brx}}}
\makeatother
\newcommand{\dd}{\mathrm{d}}
\newcommand{\E}{\mathbb{E}}
\DeclareMathOperator{\diag}{diag}
\newcommand{\tr}{\operatorname{tr}}
\newcommand{\<}{\langle}
\renewcommand{\>}{\rangle}
\usepackage{enumitem}
\usepackage{empheq}
\usepackage[utf8]{inputenc}
\usepackage[T1]{fontenc}
\usepackage{lmodern}
\usepackage{bibentry}
\usepackage{soul}
\usepackage{tikz}
\usepackage{xcolor}
\usepackage{eso-pic}
\usepackage[capitalise]{cleveref}
% \hypersetup{
%  colorlinks=true,
%  linkcolor=black,    % Color of internal links
%  citecolor=black,    % Color of citations
%  urlcolor=black,     % Color of URLs
%  anchorcolor=black,  % Color of anchors
%  pdftitle={},
%  pdfauthor={© Sven Karbach},
%  pdfsubject={},
%  pdfkeywords={}
% }
\newcommand{\abs}[1]{\left|\left|#1\right|\right|}
\newtheorem{theorem}{Theorem}[section]
\newtheorem{corollary}[theorem]{Corollary}
\newtheorem{assumption}[theorem]{Assumption}
\newtheorem{example}[theorem]{Example}
\newtheorem{lemma}[theorem]{Lemma}
\newtheorem{definition}[theorem]{Definition}
\newtheorem{remark}[theorem]{Remark}
\newtheorem{proposition}[theorem]{Proposition}
\setlength\parindent{0pt}
\newtheoremstyle{example}
  {.3\baselineskip}
  {.3\baselineskip}
  {\normalsize}  % BODYFONT
  {0pt}       % INDENT (empty value is the same as 0pt)
  {\bfseries} % HEADFONT
  {.}         % HEADPUNCT
  {5pt plus 1pt minus 1pt} % HEADSPACE
  {}          % CUSTOM-HEAD-SPEC
\theoremstyle{example}
\DeclareMathOperator\Tr{Tr}
\DeclareMathOperator\dom{dom}
\DeclareMathOperator\Res{Res}
\DeclareMathOperator\lin{lin}
% \DeclareMathOperator\diag{diag}
\DeclareMathOperator{\vect}{vec}
\newcommand*\D{\mathop{}\!\mathrm{d}}
% \newcommand*\E{\mathop{}\!\mathrm{e}}
\newcommand*\I{\mathop{}\!\mathrm{i}}
\def\e{{\mathrm{e}}}
\def\S{{\mathcal S}}
\def\R{{\mathcal R}}

\newcommand{\Cov}{\mathrm{Cov}}
\newcommand{\Id}{\mathrm{Id}}
\newcommand{\dvg}{\textbf{\mathrm{div}}}
\newcommand{\grad}{\textbf{\mathrm{grad}}}
\newcommand*\DD{\mathop{}\!\textnormal{D}}
\newcommand{\dt}{\D t }
\newcommand{\dnt}{\D^{n} t }
\newcommand{\du}{\D u\;}
\newcommand{\dx}{\D x }
\newcommand{\dnx}{\D^{n} x }
\newcommand{\dy}{\D y }
\newcommand{\dz}{\D z }
\newcommand{\dnu}{\D \nu }
\newcommand{\dlambda}{\D \lambda}
\newcommand{\dmuc}{\D \ol{\mu}}
\newcommand{\op}[1]{\mathbf{#1}}
\newcommand{\ve}{\mathbf{e}}
\newcommand{\veb}{\mathbf{\bar e}}
\newcommand{\vf}{\mathbf{f}}
\newcommand{\sven}[1]{\textcolor{red}{[Sven: #1]}}
\newcommand{\konstantinos}[1]{\textcolor{blue}{[Konstantinos: #1]}}

\begin{document}


\title[Forward-Curve Representation and Delta Hedging under the HHK Model]{Forward-Curve Representation and Delta Hedging under the HHK Model}



\address{\today{}, Korteweg-de Vries Institute for Mathematics and Informatics
  Institute at the University of Amsterdam, Science Park
, 1098 XH, Amsterdam.}      

\maketitle
\vspace{-3mm}
% \begin{center}
%   \begin{tabular}{cc}
    % \textsc{Konstantinos Chatziandreou} & \textsc{Sven Karbach} \\
%     \small[\texttt{\MakeLowercase{k.chatziandreou@uva.nl}}] &
%                                                               \small[\texttt{\MakeLowercase{sven@karbach.org}}]
%     \\ 
%   \end{tabular}
% \end{center}

\section{Forward-Curve Representation and Delta Hedging under the HHK Model}
\label{sec:hedging_hhk}

% In this section we derive the forward-curve sensitivities of the swing value under the HHK model
% \eqref{eq:model}--\eqref{eq:jumpOU} and the convex-cost payoff \eqref{eq:payoff}. 

The objective is to
obtain the exact delta of the swing value with respect to the forward curve and, from it, the natural
forward-strip hedging strategy.

\subsection{Forward prices under the HHK model}

For $0 \le t \le T$, define the forward price
\[
F_t^{[T]} := \mathbb{E}\!\left[S_T \mid \mathcal{F}_t\right].
\]
Let
\[
M_J(u) := \mathbb{E}\!\left[e^{uJ}\right]
\]
denote the moment generating function of a generic jump size $J$, whenever finite.

\begin{proposition}[HHK forward representation]
\label{prop:hhk_forward}
Under \eqref{eq:model}--\eqref{eq:jumpOU}, for every $0 \le t \le T$,
\begin{equation}
\label{eq:hhk_forward_formula}
F_t^{[T]}
=
\exp\!\Big(
f(T)
+
e^{-\alpha(T-t)} X_t
+
e^{-\beta(T-t)} Y_t
+
B(t,T)
\Big),
\end{equation}
where
\begin{equation}
\label{eq:hhk_B}
B(t,T)
:=
\frac{\sigma^2}{4\alpha}\Big(1-e^{-2\alpha(T-t)}\Big)
+
\lambda \int_0^{T-t}
\Big(
M_J(e^{-\beta u})-1
\Big)\,\D u.
\end{equation}
If, in addition, $J \sim \mathrm{Exp}(1/\mu_J)$ and $\mu_J<1$, then
\begin{equation}
\label{eq:hhk_B_exp}
B(t,T)
=
\frac{\sigma^2}{4\alpha}\Big(1-e^{-2\alpha(T-t)}\Big)
+
\frac{\lambda}{\beta}
\log\!\Bigg(
\frac{1-\mu_J e^{-\beta(T-t)}}{1-\mu_J}
\Bigg).
\end{equation}
\end{proposition}

\begin{proof}
The Ornstein--Uhlenbeck factor admits the explicit solution
\begin{equation}
\label{eq:X_solution_hedge}
X_T
=
e^{-\alpha(T-t)} X_t
+
\sigma \int_t^T e^{-\alpha(T-s)}\,\D W_s.
\end{equation}
Likewise, the spike factor satisfies
\begin{equation}
\label{eq:Y_solution_hedge}
Y_T
=
e^{-\beta(T-t)} Y_t
+
\sum_{t<\tau_\ell\le T} J_\ell e^{-\beta(T-\tau_\ell)},
\end{equation}
where $(\tau_\ell)_\ell$ are the jump times of $N$ after time $t$.

Since $S_T=\exp(f(T)+X_T+Y_T)$, conditioning on $\mathcal{F}_t$ and using the independence of the
Brownian and Poisson components yields
\begin{align}
F_t^{[T]}
&=
\mathbb{E}\!\left[
\exp\!\big(f(T)+X_T+Y_T\big)
\mid \mathcal{F}_t
\right]
\nonumber\\
&=
\exp\!\Big(
f(T)+e^{-\alpha(T-t)}X_t+e^{-\beta(T-t)}Y_t
\Big)
\nonumber\\
&\quad\times
\mathbb{E}\!\left[
\exp\!\Big(
\sigma \int_t^T e^{-\alpha(T-s)}\,\D W_s
\Big)
\right]
\,
\mathbb{E}\!\left[
\exp\!\Big(
\sum_{t<\tau_\ell\le T} J_\ell e^{-\beta(T-\tau_\ell)}
\Big)
\right].
\label{eq:forward_proof_split}
\end{align}
The Gaussian exponential moment gives
\[
\mathbb{E}\!\left[
\exp\!\Big(
\sigma \int_t^T e^{-\alpha(T-s)}\,\D W_s
\Big)
\right]
=
\exp\!\Big(
\frac{\sigma^2}{2}
\int_t^T e^{-2\alpha(T-s)}\,\D s
\Big)
=
\exp\!\Big(
\frac{\sigma^2}{4\alpha}(1-e^{-2\alpha(T-t)})
\Big).
\]
For the jump term, the exponential formula for Poisson random measures gives
\begin{align*}
\mathbb{E}\!\left[
\exp\!\Big(
\sum_{t<\tau_\ell\le T} J_\ell e^{-\beta(T-\tau_\ell)}
\Big)
\right]
&=
\exp\!\Bigg(
\lambda
\int_t^T
\Big(
M_J(e^{-\beta(T-s)})-1
\Big)\,\D s
\Bigg) \\
&=
\exp\!\Bigg(
\lambda
\int_0^{T-t}
\Big(
M_J(e^{-\beta u})-1
\Big)\,\D u
\Bigg).
\end{align*}
Substituting these two identities into \eqref{eq:forward_proof_split} yields
\eqref{eq:hhk_forward_formula}--\eqref{eq:hhk_B}.

If $J\sim \mathrm{Exp}(1/\mu_J)$, then
\[
M_J(u)=\frac{1}{1-\mu_J u},
\qquad
u<\mu_J^{-1}.
\]
Hence
\begin{align*}
\lambda \int_0^{T-t}\Big(M_J(e^{-\beta u})-1\Big)\,\D u
&=
\lambda \int_0^{T-t}
\frac{\mu_J e^{-\beta u}}{1-\mu_J e^{-\beta u}}\,\D u \\
&=
\frac{\lambda}{\beta}
\log\!\Bigg(
\frac{1-\mu_J e^{-\beta(T-t)}}{1-\mu_J}
\Bigg),
\end{align*}
which proves \eqref{eq:hhk_B_exp}.
\end{proof}

\begin{corollary}[Multiplicative forward representation]
\label{cor:mult_forward_rep}
For every $0 \le t \le T$,
\begin{equation}
\label{eq:spot_forward_ratio_general}
S_T = F_t^{[T]} \,\Lambda_{t,T},
\end{equation}
where
\begin{equation}
\label{eq:lambda_tT}
\Lambda_{t,T}
:=
\exp\!\Bigg(
\sigma \int_t^T e^{-\alpha(T-s)}\,\D W_s
+
\sum_{t<\tau_\ell\le T} J_\ell e^{-\beta(T-\tau_\ell)}
-
B(t,T)
\Bigg).
\end{equation}
In particular,
\begin{equation}
\label{eq:dS_dF_general}
\frac{\partial S_T}{\partial F_t^{[T]}}
=
\Lambda_{t,T}
=
\frac{S_T}{F_t^{[T]}}.
\end{equation}
\end{corollary}

\begin{proof}
Subtract \eqref{eq:hhk_forward_formula} from the logarithm of $S_T$, using
\eqref{eq:X_solution_hedge}--\eqref{eq:Y_solution_hedge}. This gives
\eqref{eq:spot_forward_ratio_general}. Since $S_T = F_t^{[T]} \Lambda_{t,T}$ and
$\Lambda_{t,T}$ does not depend on $F_t^{[T]}$, differentiation immediately yields
\eqref{eq:dS_dF_general}.
\end{proof}

A semimartingale form of the forward dynamics, useful for hedging, is obtained by introducing the
compensated Poisson process
\[
\widetilde{N}_t := N_t - \lambda t.
\]
Then, for fixed $T$,
\begin{equation}
\label{eq:forward_sde_hedge}
\frac{\D F_t^{[T]}}{F_{t-}^{[T]}}
=
\sigma e^{-\alpha(T-t)}\,\D W_t
+
\Big(
e^{J_t e^{-\beta(T-t)}}-1
\Big)\,\D \widetilde{N}_t.
\end{equation}
Thus the hedge instruments in our setting are naturally the forwards $F_t^{[T]}$ for the remaining
exercise maturities $T=t_i,\dots,t_{n-1}$.

\subsection{Envelope theorem and the swing value as a function of the initial forward curve}

We now write the swing value explicitly as a function of the initial forward curve on the exercise dates.
Define
\[
\mathbf{F}_0
:=
\big(
F_0^{[t_0]},F_0^{[t_1]},\dots,F_0^{[t_{n-1}]}
\big).
\]
At the exercise dates $t_0,\dots,t_{n-1}$, the admissible control set is
\[
\mathcal{Q}
:=
\Bigg\{
(q_{t_0},\dots,q_{t_{n-1}})
\,:\,
q_{t_i}\in[q_{\min},q_{\max}],
\ 
\sum_{i=0}^{n-1} q_{t_i}\le Q_{\max},
\ 
\text{all operational constraints are satisfied}
\Bigg\}.
\]
The value of the swing contract is then
\begin{equation}
\label{eq:V0_forward_curve}
V_0(\mathbf{F}_0)
=
\sup_{(q_{t_i})_{i=0}^{n-1}\in\mathcal{Q}}
\mathbb{E}\!\Bigg[
\sum_{i=0}^{n-1}
e^{-r t_i}
\Big(
q_{t_i}(S_{t_i}-K)^+
-
c q_{t_i}^{\gamma_{\mathrm{cost}}}
\Big)
\Bigg].
\end{equation}
By Corollary~\ref{cor:mult_forward_rep}, for each exercise date $t_i$ we have
\begin{equation}
\label{eq:exercise_spot_forward}
S_{t_i}
=
F_0^{[t_i]} \Lambda_{0,t_i},
\qquad
\Lambda_{0,t_i}
=
\frac{S_{t_i}}{F_0^{[t_i]}}.
\end{equation}
Hence \eqref{eq:V0_forward_curve} may be rewritten as
\begin{equation}
\label{eq:V0_forward_curve_explicit}
V_0(\mathbf{F}_0)
=
\sup_{(q_{t_i})_{i=0}^{n-1}\in\mathcal{Q}}
\mathbb{E}\!\Bigg[
\sum_{i=0}^{n-1}
e^{-r t_i}
\Big(
q_{t_i}(F_0^{[t_i]}\Lambda_{0,t_i}-K)^+
-
c q_{t_i}^{\gamma_{\mathrm{cost}}}
\Big)
\Bigg].
\end{equation}

The crucial point is that, in our model, the admissible set $\mathcal{Q}$ depends on the contract
specification (volume caps, refraction, ramping and other operational features), but does \emph{not}
depend on the initial forward curve $\mathbf{F}_0$. This is precisely the structural condition required for
the envelope-theorem argument.

\begin{theorem}[Envelope theorem]
\label{thm:envelope_general}
Let $f:\mathbb{R}^n\times\mathbb{R}^{\ell}\to\mathbb{R}$ and
$g_j:\mathbb{R}^n\times\mathbb{R}^{\ell}\to\mathbb{R}$, $j=1,\dots,m$, be continuously differentiable.
Consider the constrained optimization problem
\[
V(\alpha)
=
\sup_{x\in\mathbb{R}^n}
\big\{
f(x,\alpha): g_j(x,\alpha)\ge 0,\ j=1,\dots,m
\big\}.
\]
Let
\[
L(x,\mu,\alpha)
=
f(x,\alpha)+\sum_{j=1}^m \mu_j g_j(x,\alpha)
\]
be the associated Lagrangian. If $V$ is differentiable at $\alpha$ and $(x^*(\alpha),\mu^*(\alpha))$
is a saddle point of the Lagrangian, then for each component $\alpha_k$,
\[
\frac{\partial V}{\partial \alpha_k}(\alpha)
=
\frac{\partial L}{\partial \alpha_k}\big(x^*(\alpha),\mu^*(\alpha),\alpha\big).
\]
\end{theorem}

\subsection{Time-$0$ delta with respect to the initial forward curve}

We now apply Theorem~\ref{thm:envelope_general} to the swing problem.

\begin{proposition}[Time-$0$ forward-curve delta]
\label{prop:time0_delta}
Assume that an optimal exercise policy
\[
(q_{t_0}^*,\dots,q_{t_{n-1}}^*) \in \mathcal{Q}
\]
exists for \eqref{eq:V0_forward_curve_explicit}. Assume further that the map
\[
\mathbf{F}_0
\longmapsto
\mathbb{E}\!\Bigg[
\sum_{i=0}^{n-1}
e^{-r t_i}
\Big(
q_{t_i}^*(F_0^{[t_i]}\Lambda_{0,t_i}-K)^+
-
c (q_{t_i}^*)^{\gamma_{\mathrm{cost}}}
\Big)
\Bigg]
\]
is continuously differentiable. Then, for every $i=0,\dots,n-1$,
\begin{equation}
\label{eq:time0_delta_main}
\frac{\partial V_0}{\partial F_0^{[t_i]}}
=
e^{-r t_i}
\mathbb{E}\!\left[
q_{t_i}^*
\mathbf{1}_{\{S_{t_i}>K\}}
\frac{S_{t_i}}{F_0^{[t_i]}}
\right].
\end{equation}
Equivalently,
\begin{equation}
\label{eq:time0_delta_main_alt}
\frac{\partial V_0}{\partial F_0^{[t_i]}}
=
e^{-r t_i}
\mathbb{E}\!\left[
q_{t_i}^*
\mathbf{1}_{\{F_0^{[t_i]}\Lambda_{0,t_i}>K\}}
\Lambda_{0,t_i}
\right].
\end{equation}
\end{proposition}

\begin{proof}
Define the objective
\begin{equation}
\label{eq:f_objective_hedge}
f\big((q_{t_i})_{i=0}^{n-1},\mathbf{F}_0\big)
:=
\mathbb{E}\!\Bigg[
\sum_{i=0}^{n-1}
e^{-r t_i}
\Big(
q_{t_i}(F_0^{[t_i]}\Lambda_{0,t_i}-K)^+
-
c q_{t_i}^{\gamma_{\mathrm{cost}}}
\Big)
\Bigg].
\end{equation}
Let $g_1,\dots,g_m$ denote all the constraint functions defining $\mathcal{Q}$:
instantaneous exercise bounds, cumulative volume restriction, refraction constraints,
ramping limits and any further operational constraints. By construction,
\[
g_j\big((q_{t_i})_{i=0}^{n-1},\mathbf{F}_0\big)
=
g_j\big((q_{t_i})_{i=0}^{n-1}\big),
\qquad
j=1,\dots,m,
\]
so that
\begin{equation}
\label{eq:g_no_forward_dep}
\frac{\partial g_j}{\partial F_0^{[t_i]}} = 0,
\qquad
i=0,\dots,n-1,\quad j=1,\dots,m.
\end{equation}

Let
\[
L\big((q_{t_i})_{i=0}^{n-1},\mu,\mathbf{F}_0\big)
=
f\big((q_{t_i})_{i=0}^{n-1},\mathbf{F}_0\big)
+
\sum_{j=1}^m \mu_j g_j\big((q_{t_i})_{i=0}^{n-1}\big)
\]
be the Lagrangian. Applying Theorem~\ref{thm:envelope_general} and using
\eqref{eq:g_no_forward_dep}, we obtain
\begin{equation}
\label{eq:envelope_step_time0}
\frac{\partial V_0}{\partial F_0^{[t_i]}}
=
\frac{\partial L}{\partial F_0^{[t_i]}}
\big((q_{t_j}^*)_{j=0}^{n-1},\mu^*,\mathbf{F}_0\big)
=
\frac{\partial f}{\partial F_0^{[t_i]}}
\big((q_{t_j}^*)_{j=0}^{n-1},\mathbf{F}_0\big).
\end{equation}
Thus no derivative of the optimizer $(q_{t_i}^*)_{i=0}^{n-1}$ appears.

It remains to compute the partial derivative of $f$. Since only the $i$-th summand depends on
$F_0^{[t_i]}$,
\[
\frac{\partial f}{\partial F_0^{[t_i]}}
=
e^{-r t_i}
\frac{\partial}{\partial F_0^{[t_i]}}
\mathbb{E}\!\left[
q_{t_i}^*
(F_0^{[t_i]}\Lambda_{0,t_i}-K)^+
\right].
\]
The convex cost term has zero direct derivative with respect to $F_0^{[t_i]}$.

Because the factor $X_{t_i}$ contains a non-degenerate Gaussian component, the law of $S_{t_i}$ is
continuous, hence
\[
\mathbb{P}(S_{t_i}=K)=0.
\]
Therefore the positive-part function is differentiable almost surely at
$F_0^{[t_i]}\Lambda_{0,t_i}$, and
\[
\frac{\partial}{\partial F_0^{[t_i]}}
(F_0^{[t_i]}\Lambda_{0,t_i}-K)^+
=
\mathbf{1}_{\{F_0^{[t_i]}\Lambda_{0,t_i}>K\}}
\Lambda_{0,t_i}.
\]
Moreover,
\[
\left|
\frac{\partial}{\partial F_0^{[t_i]}}
\Big(
q_{t_i}^*
(F_0^{[t_i]}\Lambda_{0,t_i}-K)^+
\Big)
\right|
\le
q_{\max}\Lambda_{0,t_i},
\]
and since
\[
\mathbb{E}[\Lambda_{0,t_i}] = 1,
\]
the right-hand side is integrable. Hence differentiation may be interchanged with expectation by dominated
convergence. We conclude that
\begin{align*}
\frac{\partial V_0}{\partial F_0^{[t_i]}}
&=
e^{-r t_i}
\mathbb{E}\!\left[
q_{t_i}^*
\mathbf{1}_{\{F_0^{[t_i]}\Lambda_{0,t_i}>K\}}
\Lambda_{0,t_i}
\right] \\
&=
e^{-r t_i}
\mathbb{E}\!\left[
q_{t_i}^*
\mathbf{1}_{\{S_{t_i}>K\}}
\frac{S_{t_i}}{F_0^{[t_i]}}
\right],
\end{align*}
which proves \eqref{eq:time0_delta_main}--\eqref{eq:time0_delta_main_alt}.
\end{proof}

\begin{remark}
\label{rem:delta_interpretation}
Formula \eqref{eq:time0_delta_main} is the exact analogue, in our HHK setting, of the standard
envelope-theorem swing delta with respect to the initial forward curve. Relative to the no-cost,
linear-payoff case, two changes occur:
\begin{enumerate}
    \item the multiplicative factor is now the exact HHK ratio $S_{t_i}/F_0^{[t_i]}$;
    \item the positive-part payoff contributes the indicator $\mathbf{1}_{\{S_{t_i}>K\}}$.
\end{enumerate}
The convex cost term $-c q_{t_i}^{\gamma_{\mathrm{cost}}}$ does not contribute directly to the derivative with
respect to $F_0^{[t_i]}$, but it changes the optimizer $q_{t_i}^*$ and therefore affects the delta indirectly
through the optimal exercise policy.
\end{remark}

\subsection{Dynamic delta strategy}

For hedging, the relevant object is the \emph{continuation value} at a generic exercise date $t_i$.
Define
\begin{equation}
\label{eq:continuation_value}
V_i
:=
\operatorname*{ess\,sup}_{(q_{t_j})_{j=i}^{n-1}\in\mathcal{Q}_i}
\mathbb{E}\!\Bigg[
\sum_{j=i}^{n-1}
e^{-r(t_j-t_i)}
\Big(
q_{t_j}(S_{t_j}-K)^+
-
c q_{t_j}^{\gamma_{\mathrm{cost}}}
\Big)
\Bigg|
\mathcal{F}_{t_i}
\Bigg],
\end{equation}
where $\mathcal{Q}_i$ is the set of remaining admissible controls given the current contract state at
time $t_i$.

Let
\[
\mathbf{F}_i
:=
\big(
F_{t_i}^{[t_i]},F_{t_i}^{[t_{i+1}]},\dots,F_{t_i}^{[t_{n-1}]}
\big)
\]
denote the forward curve observed at time $t_i$ over the remaining exercise dates. For every $j\ge i$,
Corollary~\ref{cor:mult_forward_rep} gives
\begin{equation}
\label{eq:dynamic_forward_rep}
S_{t_j}
=
F_{t_i}^{[t_j]}\Lambda_{t_i,t_j},
\end{equation}
where
\begin{equation}
\label{eq:dynamic_lambda}
\Lambda_{t_i,t_j}
:=
\exp\!\Bigg(
\sigma \int_{t_i}^{t_j} e^{-\alpha(t_j-s)}\,\D W_s
+
\sum_{t_i<\tau_\ell\le t_j}
J_\ell e^{-\beta(t_j-\tau_\ell)}
-
B(t_i,t_j)
\Bigg).
\end{equation}

\begin{corollary}[Dynamic forward-curve delta]
\label{cor:dynamic_delta}
Assume that, for the continuation problem \eqref{eq:continuation_value}, an optimal continuation
policy exists and the corresponding value function is differentiable with respect to the remaining
forward curve $\mathbf{F}_i$. Then, for every $m=i,\dots,n-1$,
\begin{equation}
\label{eq:dynamic_delta_formula}
\Delta_i^{(m)}
:=
\frac{\partial V_i}{\partial F_{t_i}^{[t_m]}}
=
e^{-r(t_m-t_i)}
\mathbb{E}\!\left[
q_{t_m}^*
\mathbf{1}_{\{S_{t_m}>K\}}
\frac{S_{t_m}}{F_{t_i}^{[t_m]}}
\Bigg|
\mathcal{F}_{t_i}
\right].
\end{equation}
\end{corollary}

\begin{proof}
The proof is identical to that of Proposition~\ref{prop:time0_delta}, except that all quantities are now
conditional on $\mathcal{F}_{t_i}$ and the optimization is over the remaining admissible controls in
$\mathcal{Q}_i$.
\end{proof}

In particular, since $F_{t_i}^{[t_i]} = S_{t_i}$, the prompt delta is
\begin{equation}
\label{eq:prompt_delta}
\Delta_i^{(i)}
=
q_{t_i}^* \mathbf{1}_{\{S_{t_i}>K\}}.
\end{equation}

The natural delta hedge at time $t_i$ is therefore the forward strip
\begin{equation}
\label{eq:hedge_positions}
\phi_i^{(m)} := \Delta_i^{(m)},
\qquad
m=i,\dots,n-1.
\end{equation}
That is, at time $t_i$ one holds $\phi_i^{(m)}$ units of the forward maturing at $t_m$.

\subsection{Discrete-time hedging recursion}

Consider the seller of the swing contract. Let $B_{t_i}$ denote the value of the hedging portfolio
immediately \emph{after} the cash flow at time $t_i$ has been settled and the hedge has been rebalanced.
Starting from the premium
\[
B_{t_0}=V_0,
\]
and using that forwards are zero-cost instruments at inception, the self-financing recursion over one step is
\begin{equation}
\label{eq:self_financing_recursion}
B_{t_{i+1}}
=
e^{r\Delta t}
\Big(
B_{t_i}
-
\Pi_{t_i}(q_{t_i}^*)
\Big)
+
\sum_{m=i+1}^{n-1}
\phi_i^{(m)}
\Big(
F_{t_{i+1}}^{[t_m]} - F_{t_i}^{[t_m]}
\Big),
\qquad i=0,\dots,n-2.
\end{equation}
The position $\phi_i^{(m)}$ is then reset at time $t_{i+1}$ to the new delta
$\phi_{i+1}^{(m)}=\Delta_{i+1}^{(m)}$.

The local first-order hedge relation is therefore
\begin{equation}
\label{eq:first_order_hedge_relation}
V_{i+1}-\mathbb{E}[V_{i+1}\mid \mathcal{F}_{t_i}]
\approx
\sum_{m=i+1}^{n-1}
\Delta_i^{(m)}
\Big(
F_{t_{i+1}}^{[t_m]} - F_{t_i}^{[t_m]}
\Big),
\end{equation}
which identifies \eqref{eq:hedge_positions} as the exact forward-curve delta hedge.

\begin{remark}[Incomplete-market interpretation]
\label{rem:incomplete_market}
The HHK model contains jump risk, while the spot itself is not directly tradable. Consequently, the
market is incomplete and the strategy \eqref{eq:hedge_positions} is, in general, not a perfect replication
strategy. What is exact is the \emph{first-order sensitivity}: \eqref{eq:time0_delta_main} and
\eqref{eq:dynamic_delta_formula} are the exact deltas of the swing value with respect to the forward
curve in our model.
\end{remark}

\begin{remark}[Delivery-period forwards]
\label{rem:delivery_period_forwards}
In practice, electricity forwards are usually quoted for delivery periods. If
\[
F_t^{[T_1,T_2]}
=
\int_{T_1}^{T_2}
w(T;T_1,T_2)\,F_t^{[T]}\,\D T
\]
denotes the corresponding delivery-period forward, then a hedge in quoted bucket forwards is obtained
by applying the chain rule to the instantaneous-forward deltas derived above, once a parametrization of
the curve $T \mapsto F_t^{[T]}$ in terms of market quotes has been fixed.
\end{remark}
\section{Hedging with Delivery-Period Forwards Quoted in Energy Markets}
\label{sec:delivery_period_hedging}

In actual electricity markets, the hedge instruments are not the idealized instantaneous-maturity forwards
$F_t^{[T]}$, but delivery-period forwards such as daily, weekly, monthly, quarterly, or yearly baseload
products. In this section we extend the forward-curve delta calculus of
Sections~\ref{sec:hedging_hhk}--\ref{eq:first_order_hedge_relation} to such quoted market instruments.

The starting point is that a delivery-period forward is a weighted average of instantaneous forwards:
\begin{equation}
\label{eq:delivery_forward_def}
F_t^{[T_1,T_2]}
=
\int_{T_1}^{T_2}
w(T;T_1,T_2)\,F_t^{[T]}\,\D T,
\qquad 0\le t\le T_1<T_2.
\end{equation}
The weight function depends on the settlement convention. Two standard examples are:
\begin{equation}
\label{eq:delivery_weights}
w(T;T_1,T_2)
=
\frac{1}{T_2-T_1}
\quad\text{(single settlement at the end of delivery),}
\end{equation}
and
\begin{equation}
\label{eq:delivery_weights_discounted}
w(T;T_1,T_2)
=
\frac{r e^{-rT}}{e^{-rT_1}-e^{-rT_2}}
\quad\text{(continuous / instantaneous settlement).}
\end{equation}
In both cases,
\[
\int_{T_1}^{T_2} w(T;T_1,T_2)\,\D T = 1.
\]

\subsection{Delivery-period forwards under the HHK model}

Under Proposition~\ref{prop:hhk_forward}, the instantaneous forward is
\[
F_t^{[T]}
=
\exp\!\Big(
f(T)
+
e^{-\alpha(T-t)}X_t
+
e^{-\beta(T-t)}Y_t
+
B(t,T)
\Big).
\]
Hence every quoted delivery-period forward admits the exact HHK representation
\begin{equation}
\label{eq:delivery_forward_hhk}
F_t^{[T_1,T_2]}
=
\int_{T_1}^{T_2}
w(T;T_1,T_2)
\exp\!\Big(
f(T)
+
e^{-\alpha(T-t)}X_t
+
e^{-\beta(T-t)}Y_t
+
B(t,T)
\Big)\,\D T.
\end{equation}

Thus the market-quoted hedge instruments are nonlinear functionals of the HHK state $(X_t,Y_t)$, but
they are \emph{linear} functionals of the instantaneous forward curve $T \mapsto F_t^{[T]}$.

\subsection{The swing delta as a measure on delivery time}

Fix a rebalancing date $t_i$. From Corollary~\ref{cor:dynamic_delta}, the continuation value
$V_i$ has instantaneous-forward deltas
\[
\Delta_i^{(m)}
=
\frac{\partial V_i}{\partial F_{t_i}^{[t_m]}}
=
e^{-r(t_m-t_i)}
\mathbb{E}\!\left[
q_{t_m}^*
\mathbf{1}_{\{S_{t_m}>K\}}
\frac{S_{t_m}}{F_{t_i}^{[t_m]}}
\Bigg|
\mathcal{F}_{t_i}
\right],
\qquad m=i,\dots,n-1.
\]
It is convenient to collect these deltas into the atomic measure
\begin{equation}
\label{eq:delta_measure}
\mathfrak{D}_i(\D u)
:=
\sum_{m=i}^{n-1}
\Delta_i^{(m)}\,\delta_{t_m}(\D u),
\end{equation}
where $\delta_{t_m}$ denotes the Dirac mass at $t_m$.

For any deterministic perturbation $h$ of the instantaneous forward curve over the remaining exercise dates,
the first variation of the continuation value is
\begin{equation}
\label{eq:gateaux_delta_measure}
\DD V_i[h]
=
\sum_{m=i}^{n-1}
\Delta_i^{(m)}\,h(t_m)
=
\int h(u)\,\mathfrak{D}_i(\D u).
\end{equation}
Hence the swing delta is most naturally viewed as a signed measure on delivery time.

\subsection{Market-curve parametrization by quoted delivery products}

Let the hedge universe at time $t_i$ consist of $p$ quoted delivery-period forwards
\[
G_i^a := F_{t_i}^{[T_1^a,T_2^a]},
\qquad a=1,\dots,p.
\]
To express the instantaneous forward curve in terms of quoted products, we introduce deterministic
shape functions $\psi_a(\cdot)$ on the delivery horizon such that
\begin{equation}
\label{eq:curve_reconstruction}
F_{t_i}^{[u]}
=
\sum_{a=1}^p \psi_a(u)\,G_i^a,
\qquad u\in[t_i,\bar T],
\end{equation}
for some terminal date $\bar T \ge t_{n-1}$ covering all remaining exercise dates.

For consistency with quoted market prices, the shape functions must satisfy the normalization identities
\begin{equation}
\label{eq:shape_normalization}
\int_{T_1^b}^{T_2^b}
w(T;T_1^b,T_2^b)\,\psi_a(T)\,\D T
=
\delta_{ab},
\qquad a,b=1,\dots,p.
\end{equation}
Indeed, inserting \eqref{eq:curve_reconstruction} into \eqref{eq:delivery_forward_def} yields
\[
F_{t_i}^{[T_1^b,T_2^b]}
=
\sum_{a=1}^p G_i^a
\int_{T_1^b}^{T_2^b}
w(T;T_1^b,T_2^b)\,\psi_a(T)\,\D T
=
G_i^b.
\]

\begin{remark}
\label{rem:shape_choice}
The representation \eqref{eq:curve_reconstruction} is the mathematical bridge between idealized
instantaneous-forward hedging and actual market hedging. Typical choices are:
\begin{enumerate}
    \item \emph{piecewise-flat buckets:}
    \[
    \psi_a(u)=\mathbf{1}_{[T_1^a,T_2^a)}(u),
    \]
    for a partition of the delivery horizon into disjoint quoted buckets;
    \item \emph{piecewise-polynomial or spline shapes}, calibrated so that
    \eqref{eq:shape_normalization} holds;
    \item \emph{seasonality-corrected shapes}, where the intrabucket profile is prescribed
    by a deterministic seasonal pattern and only the bucket levels are quoted by the market.
\end{enumerate}
Once \eqref{eq:curve_reconstruction} is fixed, the hedge in quoted delivery forwards is fully determined.
\end{remark}

\subsection{Exact chain-rule deltas with respect to quoted delivery forwards}

We now differentiate the continuation value with respect to the quoted products $G_i^a$.

\begin{proposition}[Quoted-product delta]
\label{prop:quoted_product_delta}
Assume that at time $t_i$ the remaining instantaneous forward curve admits the representation
\eqref{eq:curve_reconstruction}. Then the continuation value $V_i$, viewed as a function of the quoted
market vector
\[
\mathbf{G}_i := (G_i^1,\dots,G_i^p),
\]
has deltas
\begin{equation}
\label{eq:quoted_delta_formula}
\Gamma_i^{(a)}
:=
\frac{\partial V_i}{\partial G_i^a}
=
\sum_{m=i}^{n-1}
\Delta_i^{(m)}\,\psi_a(t_m),
\qquad a=1,\dots,p.
\end{equation}
Equivalently, in measure form,
\begin{equation}
\label{eq:quoted_delta_measure}
\Gamma_i^{(a)}
=
\int \psi_a(u)\,\mathfrak{D}_i(\D u).
\end{equation}
\end{proposition}

\begin{proof}
By \eqref{eq:curve_reconstruction}, for each remaining exercise date $t_m$,
\[
F_{t_i}^{[t_m]}
=
\sum_{a=1}^p \psi_a(t_m)\,G_i^a.
\]
Hence, by the chain rule,
\begin{align*}
\frac{\partial V_i}{\partial G_i^a}
&=
\sum_{m=i}^{n-1}
\frac{\partial V_i}{\partial F_{t_i}^{[t_m]}}
\frac{\partial F_{t_i}^{[t_m]}}{\partial G_i^a} \\
&=
\sum_{m=i}^{n-1}
\Delta_i^{(m)}\,\psi_a(t_m),
\end{align*}
which proves \eqref{eq:quoted_delta_formula}. Formula
\eqref{eq:quoted_delta_measure} is just \eqref{eq:quoted_delta_formula} rewritten using
\eqref{eq:delta_measure}.
\end{proof}

\begin{corollary}[Piecewise-flat market buckets]
\label{cor:piecewise_flat_bucket_delta}
Suppose the quoted products form a partition of the relevant delivery horizon into disjoint buckets
\[
I_a=[T_1^a,T_2^a), \qquad a=1,\dots,p,
\]
and that the instantaneous forward curve is piecewise constant on these buckets:
\[
F_{t_i}^{[u]}
=
\sum_{a=1}^p G_i^a\,\mathbf{1}_{I_a}(u).
\]
Then
\begin{equation}
\label{eq:bucket_sum_delta}
\Gamma_i^{(a)}
=
\sum_{\substack{m=i,\dots,n-1\\ t_m \in I_a}}
\Delta_i^{(m)},
\qquad a=1,\dots,p.
\end{equation}
Thus the delta of a monthly (or quarterly) quoted forward is the sum of the exercise-date deltas falling
into that delivery bucket.
\end{corollary}

\begin{proof}
In this case, $\psi_a(u)=\mathbf{1}_{I_a}(u)$ in \eqref{eq:quoted_delta_formula}.
\end{proof}

\subsection{Dynamics of quoted delivery-period forwards}

To formulate a practical hedge and, later, a variance-optimal projection, we now derive the dynamics
of the quoted products.

Let $\mu(\D t,\D z)$ denote the jump measure associated with the marked Poisson process of jump sizes,
with compensator
\[
\lambda \nu(\D z)\,\D t,
\]
where $\nu$ is the law of a generic jump size $J$, and let
\[
\widetilde{\mu}(\D t,\D z)
:=
\mu(\D t,\D z)-\lambda\nu(\D z)\,\D t
\]
be the compensated jump measure.

For every instantaneous forward with maturity $T$, the HHK dynamics imply
\begin{equation}
\label{eq:instant_forward_prm}
\D F_t^{[T]}
=
\sigma e^{-\alpha(T-t)}F_t^{[T]}\,\D W_t
+
\int_{\mathbb{R}_+}
F_{t-}^{[T]}\Big(e^{z e^{-\beta(T-t)}}-1\Big)\,\widetilde{\mu}(\D t,\D z).
\end{equation}
Integrating over a delivery period yields the product dynamics.

\begin{proposition}[Quoted-product semimartingale dynamics]
\label{prop:quoted_product_dynamics}
For a quoted delivery-period forward
\[
G_t^a := F_t^{[T_1^a,T_2^a]},
\]
define
\begin{equation}
\label{eq:Aa_def}
A_t^a
:=
\sigma
\int_{T_1^a}^{T_2^a}
w(T;T_1^a,T_2^a)\,
e^{-\alpha(T-t)}F_t^{[T]}\,\D T,
\end{equation}
and
\begin{equation}
\label{eq:Ba_def}
B_t^a(z)
:=
\int_{T_1^a}^{T_2^a}
w(T;T_1^a,T_2^a)\,
F_{t-}^{[T]}
\Big(
e^{z e^{-\beta(T-t)}}-1
\Big)\,\D T.
\end{equation}
Then
\begin{equation}
\label{eq:quoted_product_semimart}
\D G_t^a
=
A_t^a\,\D W_t
+
\int_{\mathbb{R}_+}
B_t^a(z)\,\widetilde{\mu}(\D t,\D z).
\end{equation}
Consequently, the predictable quadratic covariation of two quoted products is
\begin{equation}
\label{eq:quoted_product_pqc}
\D\langle G^a,G^b\rangle_t^{p}
=
A_t^aA_t^b\,\D t
+
\lambda
\int_{\mathbb{R}_+}
B_t^a(z)B_t^b(z)\,\nu(\D z)\,\D t.
\end{equation}
\end{proposition}

\begin{proof}
Insert \eqref{eq:instant_forward_prm} into \eqref{eq:delivery_forward_def} and apply stochastic Fubini's
theorem. Formula \eqref{eq:quoted_product_pqc} follows from the orthogonality of the Brownian and
compensated-jump martingale parts.
\end{proof}

For later use, the predictable quadratic covariation between a quoted product and an instantaneous forward
with maturity $u \ge t$ is
\begin{equation}
\label{eq:cross_pqc_product_inst}
\D\langle G^a,F^{[u]}\rangle_t^{p}
=
A_t^a\,
\sigma e^{-\alpha(u-t)}F_t^{[u]}\,\D t
+
\lambda
\int_{\mathbb{R}_+}
B_t^a(z)\,
F_{t-}^{[u]}
\Big(
e^{z e^{-\beta(u-t)}}-1
\Big)\,
\nu(\D z)\,\D t.
\end{equation}

\subsection{Exact first-order hedge in quoted products}

Let
\[
\boldsymbol{\Gamma}_i
:=
\big(\Gamma_i^{(1)},\dots,\Gamma_i^{(p)}\big)^\top
\]
be the vector of quoted-product deltas from Proposition~\ref{prop:quoted_product_delta}.
Define the hedge positions
\begin{equation}
\label{eq:quoted_hedge_positions}
\eta_i^{(a)} := \Gamma_i^{(a)},
\qquad a=1,\dots,p.
\end{equation}
Then, for every admissible curve perturbation generated by variations of the quoted market vector
$\mathbf{G}_i$, the first-order change of the continuation value satisfies
\begin{equation}
\label{eq:exact_first_order_quoted}
\DD V_i[h_{\delta \mathbf{G}}]
=
\sum_{a=1}^p \eta_i^{(a)}\,\delta G_i^a,
\end{equation}
where
\[
h_{\delta \mathbf{G}}(u)
=
\sum_{a=1}^p \psi_a(u)\,\delta G_i^a.
\]
Thus \eqref{eq:quoted_hedge_positions} is the exact first-order hedge \emph{within the chosen market
representation of the forward curve}.

In discrete time, the self-financing hedge recursion becomes
\begin{equation}
\label{eq:self_financing_quoted}
B_{t_{i+1}}
=
e^{r\Delta t}
\Big(
B_{t_i}-\Pi_{t_i}(q_{t_i}^*)
\Big)
+
\sum_{a=1}^p
\eta_i^{(a)}
\Big(
G_{t_{i+1}}^a-G_{t_i}^a
\Big),
\qquad i=0,\dots,n-2.
\end{equation}

\subsection{Variance-optimal hedge when the quoted products do not span the full curve}

In practice, the quoted hedge universe may be too small to describe all exercise-date forward exposures
exactly. Then one should project the idealized instantaneous-forward hedge onto the span of the quoted
products.

Let
\[
\Delta G_i^a := G_{t_{i+1}}^a-G_{t_i}^a,
\qquad
\Delta \mathbf{G}_i := (\Delta G_i^1,\dots,\Delta G_i^p)^\top,
\]
and define the centered one-step change of the continuation value by
\[
\Delta \widetilde V_i
:=
V_{i+1}-\mathbb{E}[V_{i+1}\mid \mathcal{F}_{t_i}].
\]
The locally variance-optimal quoted-product hedge is the solution of
\begin{equation}
\label{eq:vo_problem}
\eta_i^{\mathrm{vo}}
=
\operatorname*{arg\,min}_{\eta\in\mathbb{R}^p}
\mathbb{E}\!\left[
\big(
\Delta \widetilde V_i-\eta^\top \Delta \mathbf{G}_i
\big)^2
\Big|
\mathcal{F}_{t_i}
\right].
\end{equation}

\begin{proposition}[Local variance-optimal quoted-product hedge]
\label{prop:vo_quoted}
Assume the conditional covariance matrix
\[
C_i
:=
\mathbb{E}\!\left[
\Delta \mathbf{G}_i\,\Delta \mathbf{G}_i^\top
\Big|
\mathcal{F}_{t_i}
\right]
\]
is invertible. Then the unique solution to \eqref{eq:vo_problem} is
\begin{equation}
\label{eq:vo_solution}
\eta_i^{\mathrm{vo}}
=
C_i^{-1} b_i,
\qquad
b_i
:=
\mathbb{E}\!\left[
\Delta \mathbf{G}_i\,\Delta \widetilde V_i
\Big|
\mathcal{F}_{t_i}
\right].
\end{equation}
\end{proposition}

\begin{proof}
Expand the conditional quadratic criterion in \eqref{eq:vo_problem} and differentiate with respect to
$\eta$. The normal equations are
\[
C_i \eta = b_i.
\]
Since $C_i$ is invertible, the solution is \eqref{eq:vo_solution}.
\end{proof}

To connect \eqref{eq:vo_solution} with the delta formulas derived earlier, use the first-order
approximation
\begin{equation}
\label{eq:delta_linearization_value}
\Delta \widetilde V_i
\approx
\sum_{m=i}^{n-1}
\Delta_i^{(m)}
\Big(
F_{t_{i+1}}^{[t_m]}-F_{t_i}^{[t_m]}
\Big).
\end{equation}
Then
\begin{equation}
\label{eq:bi_delta_approx}
b_i^a
\approx
\sum_{m=i}^{n-1}
\Delta_i^{(m)}
\,
\mathbb{E}\!\left[
\Delta G_i^a
\Big(
F_{t_{i+1}}^{[t_m]}-F_{t_i}^{[t_m]}
\Big)
\Big|
\mathcal{F}_{t_i}
\right].
\end{equation}
The conditional covariances entering \eqref{eq:bi_delta_approx} are determined, to first order in
$\Delta t$, by the predictable covariations
\eqref{eq:quoted_product_pqc} and \eqref{eq:cross_pqc_product_inst}.

\begin{corollary}[Reduction to the exact chain-rule hedge]
\label{cor:vo_equals_chain_rule}
Suppose the market-curve parametrization \eqref{eq:curve_reconstruction} is exact over the relevant
delivery horizon, so that for one-step changes
\[
F_{t_{i+1}}^{[t_m]}-F_{t_i}^{[t_m]}
=
\sum_{a=1}^p
\psi_a(t_m)\,
\Delta G_i^a,
\qquad m=i,\dots,n-1.
\]
Then
\[
\eta_i^{\mathrm{vo}}=\boldsymbol{\Gamma}_i.
\]
Hence the local variance-optimal hedge coincides with the exact chain-rule quoted-product delta hedge.
\end{corollary}

\begin{proof}
Using the exact representation of the forward increments and \eqref{eq:delta_linearization_value},
\[
\Delta \widetilde V_i
\approx
\sum_{m=i}^{n-1}
\Delta_i^{(m)}
\sum_{a=1}^p \psi_a(t_m)\Delta G_i^a
=
\sum_{a=1}^p
\Gamma_i^{(a)}\Delta G_i^a
=
\boldsymbol{\Gamma}_i^\top \Delta \mathbf{G}_i.
\]
Therefore,
\[
b_i
=
\mathbb{E}\!\left[
\Delta \mathbf{G}_i\Delta \mathbf{G}_i^\top
\Big|
\mathcal{F}_{t_i}
\right]
\boldsymbol{\Gamma}_i
=
C_i \boldsymbol{\Gamma}_i,
\]
and \eqref{eq:vo_solution} yields
\[
\eta_i^{\mathrm{vo}}=\boldsymbol{\Gamma}_i.
\]
\end{proof}

\subsection{Interpretation}

The extension from instantaneous forwards to actual market instruments is therefore mathematically clean.

First, one computes the ideal swing deltas
\[
\Delta_i^{(m)} = \frac{\partial V_i}{\partial F_{t_i}^{[t_m]}}
\]
at the exercise dates.

Second, one specifies how the market quotes $(G_i^a)_{a=1}^p$ reconstruct the instantaneous forward
curve through \eqref{eq:curve_reconstruction}. This induces the quoted-product deltas
\[
\Gamma_i^{(a)}
=
\sum_{m=i}^{n-1}
\Delta_i^{(m)}\psi_a(t_m).
\]

Third, one hedges in the actually traded delivery-period forwards using either:
\begin{enumerate}
    \item the exact first-order chain-rule hedge $\eta_i^{(a)}=\Gamma_i^{(a)}$ when the chosen market
    parametrization spans the relevant forward-curve deformations, or
    \item the local variance-optimal projection $\eta_i^{\mathrm{vo}}=C_i^{-1}b_i$ when the hedge
    universe is incomplete.
\end{enumerate}

Thus the mathematically correct object being hedged is not a single scalar delta, but the forward-curve
delta measure \eqref{eq:delta_measure}; the market hedge is obtained by pushing this measure through
the quoted-product basis of the delivery curve.


%========================================================
% INSERT THIS AFTER THE DELIVERY-PERIOD HEDGING SECTION
%========================================================

\subsection{RL-Based Implementation of Pricing and Hedging}
\label{sec:rl_pricing_hedging_impl}

We now give an implementable algorithmic procedure that connects the pricing policy learned by deep
reinforcement learning to the forward-curve hedge derived in
Sections~\ref{sec:hedging_hhk}--\ref{sec:delivery_period_hedging}. The key point is that the trained
actor $\pi_\theta$ is used in two roles:

\begin{enumerate}
    \item as the policy generating the optimal exercise quantities along simulated HHK paths, and
    \item as the policy defining the continuation-value and delta estimators used to construct the hedge.
\end{enumerate}

At each exercise date $t_i$, let $s_i$ denote the pre-exercise state and let
\[
q_{t_i}^{\pi_\theta} := \mathcal{A}(s_i,\pi_\theta(s_i))
\]
be the feasible action produced by the trained actor after denormalization, profitability gating,
cumulative-volume clipping, and operational masking. Here $\mathcal{A}$ denotes the full feasibility map
implemented by the environment and the actor gate from Section~\ref{sec:profitability_gate}.

For a quoted hedge universe
\[
\mathbf{G}_i = (G_i^1,\dots,G_i^p)^\top,
\qquad
G_i^a = F_{t_i}^{[T_1^a,T_2^a]},
\]
assume the remaining instantaneous forward curve is reconstructed through
\begin{equation}
\label{eq:rl_impl_curve_reconstruction}
F_{t_i}^{[u]}
=
\sum_{a=1}^p \psi_a(u)\,G_i^a,
\qquad
u\in[t_i,\bar T].
\end{equation}

\paragraph{Policy rollout price estimator.}
For any state $s_i$ and quoted forward vector $\mathbf{G}_i$, define the risk-neutral rollout price under
the learned policy by
\begin{equation}
\label{eq:policy_value_estimator}
\widehat{V}_i^{\pi_\theta}(s_i,\mathbf{G}_i)
:=
\frac{1}{M_V}
\sum_{\ell=1}^{M_V}
\sum_{m=i}^{n-1}
e^{-r(t_m-t_i)}
\Big(
q_{t_m}^{(\ell)}(S_{t_m}^{(\ell)}-K)^+
-
c \big(q_{t_m}^{(\ell)}\big)^{\gamma_{\mathrm{cost}}}
\Big),
\end{equation}
where, on rollout $\ell$, the process
\[
\big(X_{t_m}^{(\ell)},Y_{t_m}^{(\ell)},S_{t_m}^{(\ell)}\big)_{m=i}^{n-1}
\]
is simulated under the HHK model conditional on the current state at time $t_i$, and the actions
$q_{t_m}^{(\ell)}$ are generated recursively by the feasible actor policy.

\paragraph{Continuation-value estimator after the current exercise.}
At time $t_i$, once the current exercise $q_{t_i}^{\pi_\theta}$ has been fixed and the inventory / refraction
state has been updated, let $\bar s_i$ denote the post-exercise state. The continuation value is estimated by
\begin{equation}
\label{eq:continuation_value_estimator_impl}
\widehat{V}_{i,\mathrm{cont}}^{\pi_\theta}(\bar s_i,\mathbf{G}_i)
:=
\frac{1}{M_\Delta}
\sum_{\ell=1}^{M_\Delta}
\sum_{m=i+1}^{n-1}
e^{-r(t_m-t_i)}
\Big(
q_{t_m}^{(\ell)}(S_{t_m}^{(\ell)}-K)^+
-
c \big(q_{t_m}^{(\ell)}\big)^{\gamma_{\mathrm{cost}}}
\Big).
\end{equation}

\paragraph{Instantaneous-forward delta estimator.}
Using Corollary~\ref{cor:dynamic_delta}, for every future exercise date $t_m$, $m=i+1,\dots,n-1$, define
the Monte Carlo estimator
\begin{equation}
\label{eq:delta_mc_estimator_impl}
\widehat{\Delta}_i^{(m)}
:=
\frac{e^{-r(t_m-t_i)}}{M_\Delta\,F_{t_i}^{[t_m]}}
\sum_{\ell=1}^{M_\Delta}
q_{t_m}^{(\ell)}
\mathbf{1}_{\{S_{t_m}^{(\ell)}>K\}}
S_{t_m}^{(\ell)}.
\end{equation}
This is the discrete implementation of
\[
\Delta_i^{(m)}
=
e^{-r(t_m-t_i)}
\mathbb{E}\!\left[
q_{t_m}^*
\mathbf{1}_{\{S_{t_m}>K\}}
\frac{S_{t_m}}{F_{t_i}^{[t_m]}}
\Bigg|
\mathcal{F}_{t_i}
\right].
\]

\paragraph{Quoted-product delta estimator.}
Define the shape matrix
\begin{equation}
\label{eq:shape_matrix_impl}
\Psi_i
\in \mathbb{R}^{p\times(n-i-1)},
\qquad
[\Psi_i]_{a,m-i}
:=
\psi_a(t_m),
\qquad
a=1,\dots,p,\quad m=i+1,\dots,n-1.
\end{equation}
Let
\[
\widehat{\boldsymbol{\Delta}}_i
:=
\big(
\widehat{\Delta}_i^{(i+1)},\dots,\widehat{\Delta}_i^{(n-1)}
\big)^\top.
\]
Then the exact first-order quoted-product hedge is
\begin{equation}
\label{eq:quoted_delta_mc_impl}
\widehat{\boldsymbol{\Gamma}}_i
=
\Psi_i \widehat{\boldsymbol{\Delta}}_i,
\qquad
\widehat{\Gamma}_i^{(a)}
=
\sum_{m=i+1}^{n-1}
\widehat{\Delta}_i^{(m)}\psi_a(t_m).
\end{equation}

\paragraph{Variance-optimal projection.}
If the quoted hedge universe does not span the relevant forward-curve deformations, then the hedge is
obtained by the local variance-optimal projection. For one-step scenarios
$\ell=1,\dots,M_{\mathrm{vo}}$, define
\[
\Delta \mathbf{G}_i^{(\ell)}
:=
\mathbf{G}_{i+1}^{(\ell)}-\mathbf{G}_i,
\qquad
\Delta \widetilde V_i^{(\ell)}
:=
\widehat{V}_{i+1}^{\pi_\theta}\!\big(s_{i+1}^{(\ell)},\mathbf{G}_{i+1}^{(\ell)}\big)
-
\frac{1}{M_{\mathrm{vo}}}
\sum_{k=1}^{M_{\mathrm{vo}}}
\widehat{V}_{i+1}^{\pi_\theta}\!\big(s_{i+1}^{(k)},\mathbf{G}_{i+1}^{(k)}\big).
\]
Then
\begin{equation}
\label{eq:vo_cov_est_impl}
\widehat{C}_i
:=
\frac{1}{M_{\mathrm{vo}}-1}
\sum_{\ell=1}^{M_{\mathrm{vo}}}
\big(\Delta \mathbf{G}_i^{(\ell)}-\overline{\Delta \mathbf{G}}_i\big)
\big(\Delta \mathbf{G}_i^{(\ell)}-\overline{\Delta \mathbf{G}}_i\big)^\top,
\end{equation}
\begin{equation}
\label{eq:vo_cross_est_impl}
\widehat{b}_i
:=
\frac{1}{M_{\mathrm{vo}}-1}
\sum_{\ell=1}^{M_{\mathrm{vo}}}
\big(\Delta \mathbf{G}_i^{(\ell)}-\overline{\Delta \mathbf{G}}_i\big)
\big(\Delta \widetilde V_i^{(\ell)}-\overline{\Delta \widetilde V}_i\big),
\end{equation}
and the regularized variance-optimal hedge is
\begin{equation}
\label{eq:vo_hedge_impl}
\widehat{\boldsymbol{\eta}}_i^{\mathrm{vo}}
=
\big(\widehat{C}_i+\varepsilon_{\mathrm{reg}} I_p\big)^{-1}\widehat{b}_i.
\end{equation}

\paragraph{Self-financing update.}
Let $B_{t_i}$ denote the hedging portfolio value at time $t_i$, and let
$\boldsymbol{\eta}_i$ be the quoted-product hedge chosen at time $t_i$. Then the discrete-time
self-financing recursion is
\begin{equation}
\label{eq:self_financing_impl}
B_{t_{i+1}}
=
e^{r\Delta t}
\Big(
B_{t_i} - \Pi_{t_i}(q_{t_i}^{\pi_\theta})
\Big)
+
\boldsymbol{\eta}_i^\top
\big(
\mathbf{G}_{i+1}-\mathbf{G}_i
\big),
\qquad i=0,\dots,n-2.
\end{equation}

\paragraph{Hedge error.}
For diagnostics, define the post-rebalancing mark-to-model hedge error by
\begin{equation}
\label{eq:hedge_error_impl}
\varepsilon_i^{\mathrm{hedge}}
:=
B_{t_i}
-
\widehat{V}_i^{\pi_\theta}(s_i,\mathbf{G}_i).
\end{equation}
\begin{comment}
Algorithm~\ref{alg:rl_training_pricing} trains the pricing policy.
Algorithm~\ref{alg:cont_value_delta_estimation} estimates continuation values and deltas from the trained
policy. Algorithm~\ref{alg:rl_online_hedging} gives the full online pricing-and-hedging procedure.

\begin{algorithm}[htbp]
\caption{Offline RL Training and Policy Selection for Pricing}
\label{alg:rl_training_pricing}
\begin{algorithmic}[1]
\Require HHK parameters $\Theta_{\mathrm{HHK}}$, contract parameters $\mathcal{C}$, training episodes $N_{\mathrm{train}}$, evaluation paths $N_{\mathrm{eval}}$
\Require Replay parameters, optimizer parameters, profitability-gate parameters, shape functions $\{\psi_a\}_{a=1}^p$
\Ensure Trained actor $\pi_{\theta^\star}$, critic $Q_{\phi^\star}$, preprocessing statistics, selected checkpoint $(\theta^\star,\phi^\star)$
\State Generate risk-neutral HHK training and evaluation paths
\State Initialize actor $\pi_\theta$, critic $Q_\phi$, target networks $\pi_{\theta'}$, $Q_{\phi'}$
\State Initialize replay buffer $\mathcal{D}$ and input preprocessors
\State Set actor feasibility parameters $(q_{\min},q_{\max},Q_{\max},c,\gamma_{\mathrm{cost}})$
\For{episode $e=1$ to $N_{\mathrm{train}}$}
    \State Reset environment to initial state $s_0^{(e)}$
    \For{decision step $i=0$ to $n-1$}
        \State Compute actor output $\tilde q_{t_i} \gets \pi_\theta(s_i)$ with pre-squash exploration noise
        \State Map to feasible action $q_{t_i} \gets \mathcal{A}(s_i,\tilde q_{t_i})$
        \State Execute $q_{t_i}$ in the environment and observe $(s_{i+1},r_i,\mathrm{done})$
        \State Store transition $(s_i,q_{t_i},r_i,s_{i+1},\mathrm{done})$ in $\mathcal{D}$
        \If{replay buffer size is at least $N_{\min}$ and update condition is satisfied}
            \State Sample a minibatch from $\mathcal{D}$ with importance weights
            \State Update critic parameters $\phi$ by minimizing the TD loss
            \If{$e > e_{\mathrm{warmup}}$}
                \State Update actor parameters $\theta$ by deterministic policy gradient
                \State Soft-update actor target $\theta' \gets (1-\tau)\theta' + \tau\theta$
            \EndIf
            \State Soft-update critic target $\phi' \gets (1-\tau)\phi' + \tau\phi$
        \EndIf
    \EndFor
    \If{$e$ is an evaluation checkpoint}
        \State Estimate out-of-sample price $\widehat V_0^{(e)}$ by policy rollouts on evaluation paths
        \State Save checkpoint $(\theta_e,\phi_e,\widehat V_0^{(e)})$
    \EndIf
\EndFor
\State Select best checkpoint
\[
(\theta^\star,\phi^\star)
\gets
\arg\max_{e\in\mathcal E_{\mathrm{chk}}} \widehat V_0^{(e)}
\]
\State \Return $\pi_{\theta^\star}$, $Q_{\phi^\star}$, preprocessing statistics, selected checkpoint
\end{algorithmic}
\end{algorithm}

\begin{algorithm}[htbp]
\caption{Continuation-Value and Delta Estimation from the Trained RL Policy}
\label{alg:cont_value_delta_estimation}
\begin{algorithmic}[1]
\Require Post-exercise state $\bar s_i$ at time $t_i$
\Require Current quoted forward vector $\mathbf{G}_i=(G_i^1,\dots,G_i^p)^\top$
\Require Trained actor $\pi_\theta$, shape functions $\{\psi_a\}_{a=1}^p$, rollout count $M_\Delta$
\Ensure Continuation-value estimate $\widehat V_{i,\mathrm{cont}}^{\pi_\theta}$, instantaneous-forward deltas $\widehat{\boldsymbol{\Delta}}_i$, quoted-product deltas $\widehat{\boldsymbol{\Gamma}}_i$
\For{$m=i+1$ to $n-1$}
    \State Reconstruct instantaneous forward at exercise date $t_m$:
    \[
    F_{t_i}^{[t_m]} \gets \sum_{a=1}^{p} \psi_a(t_m)\,G_i^a
    \]
    \State Initialize accumulators $A_m \gets 0$
\EndFor
\State Initialize continuation-value accumulator $V_{\mathrm{sum}} \gets 0$
\For{rollout $\ell=1$ to $M_\Delta$}
    \State Initialize simulated HHK state at time $t_i$ from $\bar s_i$
    \State Initialize a copy of the post-exercise contract state
    \For{$m=i+1$ to $n-1$}
        \State Simulate $(X_{t_m}^{(\ell)},Y_{t_m}^{(\ell)},S_{t_m}^{(\ell)})$ under HHK from $t_{m-1}$ to $t_m$
        \State Construct RL state $s_{t_m}^{(\ell)}$ from the simulated factors and current contract variables
        \State Compute feasible actor action
        \[
        q_{t_m}^{(\ell)} \gets \mathcal{A}\big(s_{t_m}^{(\ell)},\pi_\theta(s_{t_m}^{(\ell)})\big)
        \]
        \State Update continuation-value accumulator
        \[
        V_{\mathrm{sum}}
        \gets
        V_{\mathrm{sum}}
        +
        e^{-r(t_m-t_i)}
        \Big(
        q_{t_m}^{(\ell)}(S_{t_m}^{(\ell)}-K)^+
        -
        c\big(q_{t_m}^{(\ell)}\big)^{\gamma_{\mathrm{cost}}}
        \Big)
        \]
        \State Update delta numerator accumulator
        \[
        A_m
        \gets
        A_m
        +
        e^{-r(t_m-t_i)}
        q_{t_m}^{(\ell)}
        \mathbf{1}_{\{S_{t_m}^{(\ell)}>K\}}
        S_{t_m}^{(\ell)}
        \]
        \State Update the copied contract variables using $q_{t_m}^{(\ell)}$
    \EndFor
\EndFor
\State Compute continuation value
\[
\widehat V_{i,\mathrm{cont}}^{\pi_\theta}
\gets
\frac{V_{\mathrm{sum}}}{M_\Delta}
\]
\For{$m=i+1$ to $n-1$}
    \State Compute instantaneous-forward delta
    \[
    \widehat{\Delta}_i^{(m)}
    \gets
    \frac{A_m}{M_\Delta\,F_{t_i}^{[t_m]}}
    \]
\EndFor
\State Form delta vector
\[
\widehat{\boldsymbol{\Delta}}_i
\gets
\big(
\widehat{\Delta}_i^{(i+1)},\dots,\widehat{\Delta}_i^{(n-1)}
\big)^\top
\]
\State Build shape matrix $\Psi_i$ with entries $[\Psi_i]_{a,m-i}=\psi_a(t_m)$
\State Compute quoted-product deltas
\[
\widehat{\boldsymbol{\Gamma}}_i \gets \Psi_i \widehat{\boldsymbol{\Delta}}_i
\]
\State \Return $\widehat V_{i,\mathrm{cont}}^{\pi_\theta}$, $\widehat{\boldsymbol{\Delta}}_i$, $\widehat{\boldsymbol{\Gamma}}_i$
\end{algorithmic}
\end{algorithm}

\begin{algorithm}[htbp]
\caption{Online RL Pricing and Hedge Rebalancing in Quoted Delivery-Period Forwards}
\label{alg:rl_online_hedging}
\begin{algorithmic}[1]
\Require Initial state $s_0$, initial quoted forward vector $\mathbf{G}_0$
\Require Trained actor $\pi_\theta$, shape functions $\{\psi_a\}_{a=1}^p$
\Require Pricing rollout count $M_V$, delta rollout count $M_\Delta$, variance-optimal rollout count $M_{\mathrm{vo}}$
\Require Hedge mode $\mathrm{mode}\in\{\mathrm{exact},\mathrm{vo}\}$ and regularization level $\varepsilon_{\mathrm{reg}}>0$
\Ensure Initial price $\widehat V_0^{\pi_\theta}$, hedge positions $\{\boldsymbol{\eta}_i\}_{i=0}^{n-2}$, bank-account path $\{B_{t_i}\}$, hedge errors $\{\varepsilon_i^{\mathrm{hedge}}\}$
\State Compute initial price by policy rollouts
\[
\widehat V_0^{\pi_\theta}
\gets
\widehat V_0^{\pi_\theta}(s_0,\mathbf{G}_0)
\]
\State Initialize bank account $B_{t_0} \gets \widehat V_0^{\pi_\theta}$
\For{$i=0$ to $n-2$}
    \State Observe current state $s_i$ and quoted forward vector $\mathbf{G}_i$
    \State Compute feasible current exercise
    \[
    q_{t_i}^{\pi_\theta}
    \gets
    \mathcal{A}(s_i,\pi_\theta(s_i))
    \]
    \State Compute current cash flow
    \[
    C_i
    \gets
    \Pi_{t_i}(q_{t_i}^{\pi_\theta})
    \]
    \State Update to post-exercise state $\bar s_i$
    \State Estimate continuation value and exact quoted-product deltas by Algorithm~\ref{alg:cont_value_delta_estimation}
    \Statex \hspace{\algorithmicindent} $(\widehat V_{i,\mathrm{cont}}^{\pi_\theta},\widehat{\boldsymbol{\Delta}}_i,\widehat{\boldsymbol{\Gamma}}_i)
    \gets
    \textsc{ContinuationAndDeltas}(\bar s_i,\mathbf{G}_i,\pi_\theta,\{\psi_a\},M_\Delta)$
    \If{$\mathrm{mode}=\mathrm{exact}$}
        \State Set hedge
        \[
        \boldsymbol{\eta}_i \gets \widehat{\boldsymbol{\Gamma}}_i
        \]
    \Else
        \For{scenario $\ell=1$ to $M_{\mathrm{vo}}$}
            \State Simulate one-step HHK evolution from $\bar s_i$ to obtain next state $s_{i+1}^{(\ell)}$
            \State Simulate / reconstruct next quoted forward vector $\mathbf{G}_{i+1}^{(\ell)}$
            \State Compute quoted-product increment
            \[
            \Delta \mathbf{G}_i^{(\ell)}
            \gets
            \mathbf{G}_{i+1}^{(\ell)}-\mathbf{G}_i
            \]
            \State Estimate next continuation value under the same actor policy
            \[
            U_\ell
            \gets
            \widehat V_{i+1}^{\pi_\theta}\big(s_{i+1}^{(\ell)},\mathbf{G}_{i+1}^{(\ell)}\big)
            \]
        \EndFor
        \State Compute sample means
        \[
        \overline{\Delta \mathbf{G}}_i
        \gets
        \frac{1}{M_{\mathrm{vo}}}\sum_{\ell=1}^{M_{\mathrm{vo}}}\Delta \mathbf{G}_i^{(\ell)},
        \qquad
        \overline U
        \gets
        \frac{1}{M_{\mathrm{vo}}}\sum_{\ell=1}^{M_{\mathrm{vo}}} U_\ell
        \]
        \State Initialize $\widehat C_i \gets 0_{p\times p}$ and $\widehat b_i \gets 0_p$
        \For{scenario $\ell=1$ to $M_{\mathrm{vo}}$}
            \State Center increments
            \[
            \widetilde{\Delta \mathbf{G}}_i^{(\ell)}
            \gets
            \Delta \mathbf{G}_i^{(\ell)}-\overline{\Delta \mathbf{G}}_i,
            \qquad
            \widetilde U_\ell
            \gets
            U_\ell-\overline U
            \]
            \State Accumulate covariance terms
            \[
            \widehat C_i
            \gets
            \widehat C_i
            +
            \widetilde{\Delta \mathbf{G}}_i^{(\ell)}
            \big(\widetilde{\Delta \mathbf{G}}_i^{(\ell)}\big)^\top
            \]
            \[
            \widehat b_i
            \gets
            \widehat b_i
            +
            \widetilde{\Delta \mathbf{G}}_i^{(\ell)}\,\widetilde U_\ell
            \]
        \EndFor
        \State Normalize
        \[
        \widehat C_i \gets \frac{1}{M_{\mathrm{vo}}-1}\widehat C_i,
        \qquad
        \widehat b_i \gets \frac{1}{M_{\mathrm{vo}}-1}\widehat b_i
        \]
        \State Solve the regularized normal equations
        \[
        \boldsymbol{\eta}_i
        \gets
        \big(\widehat C_i+\varepsilon_{\mathrm{reg}}I_p\big)^{-1}\widehat b_i
        \]
    \EndIf
    \State Observe / simulate realized next quoted forward vector $\mathbf{G}_{i+1}$ and realized next state $s_{i+1}$
    \State Update bank account by the self-financing recursion
    \[
    B_{t_{i+1}}
    \gets
    e^{r\Delta t}\big(B_{t_i}-C_i\big)
    +
    \boldsymbol{\eta}_i^\top(\mathbf{G}_{i+1}-\mathbf{G}_i)
    \]
    \State Compute hedge error diagnostic
    \[
    \varepsilon_{i+1}^{\mathrm{hedge}}
    \gets
    B_{t_{i+1}}
    -
    \widehat V_{i+1}^{\pi_\theta}(s_{i+1},\mathbf{G}_{i+1})
    \]
\EndFor
\State At the final exercise date $t_{n-1}$, compute and settle
\[
q_{t_{n-1}}^{\pi_\theta} \gets \mathcal{A}(s_{n-1},\pi_\theta(s_{n-1})),
\qquad
C_{n-1} \gets \Pi_{t_{n-1}}(q_{t_{n-1}}^{\pi_\theta})
\]
\State Final bank account:
\[
B_{t_n}
\gets
e^{r\Delta t}\big(B_{t_{n-1}}-C_{n-1}\big)
\]
\State \Return $\widehat V_0^{\pi_\theta}$, $\{\boldsymbol{\eta}_i\}_{i=0}^{n-2}$, $\{B_{t_i}\}_{i=0}^{n}$, $\{\varepsilon_i^{\mathrm{hedge}}\}_{i=1}^{n-1}$
\end{algorithmic}
\end{algorithm}

\paragraph{Implementation remarks.}
Algorithm~\ref{alg:rl_online_hedging} is the mathematically faithful implementation of the theory.
The exact-mode branch uses the chain-rule hedge
\[
\boldsymbol{\eta}_i = \widehat{\boldsymbol{\Gamma}}_i,
\]
while the $\mathrm{vo}$-branch uses the local variance-optimal projection
\[
\boldsymbol{\eta}_i
=
\big(\widehat C_i+\varepsilon_{\mathrm{reg}}I_p\big)^{-1}\widehat b_i.
\]
If computational speed is critical, the rollout estimator
$\widehat V_i^{\pi_\theta}(s_i,\mathbf{G}_i)$ may be replaced by the critic approximation
\[
Q_\phi\!\big(s_i,\mathcal{A}(s_i,\pi_\theta(s_i))\big),
\]
but the rollout version is the cleaner implementation because it remains fully aligned with the
risk-neutral control problem used for pricing.

\end{comment}
\end{document}